%%%%%%%%%%%%Outline%%%%%%%%%%%%%%%%%%

%%%
% message<\tool recap: dual-form operations under the existing verifier, measured wins, modules>
% we presented \tool, extending the pipeline with operations the kernel JIT emits directly
% each operation carries a verifier-checked proof sequence and a native emit the CPU runs
% only the native emit joins the trusted computing base
% with \tool we built \lang, seven hardware-idiom operations
% Lean 4 proofs show each \lang emit computes the same result as its proof sequence
% \lang speeds up eBPF programs up to 24%/22%, recovering up to 42% of the gap
% a new operation ships as a kernel module, the settled safety argument stays untouched

%%%%%%%%%%%%Outline%%%%%%%%%%%%%%%%%%

\section{Conclusion}
\label{sec:conclusion}

We presented \tool, a mechanism that extends the eBPF compilation pipeline with new operations the kernel JIT emits directly. Each operation carries a proof sequence that the existing verifier checks on every load and a native emit that the CPU runs. Beyond the generic enforcement core, the native emit is the only per-operation addition to the trusted base. With \tool, we built \lang, a set of seven hardware-idiom operations. Lean 4 proofs show that each \lang native emit has the same eBPF-visible effect as its proof sequence. \lang speeds up eBPF microbenchmarks by up to 24\% on x86-64 and 22\% on ARM64 over the unmodified pipeline, recovering up to 42\% of the gap to native code, and improves production application throughput by up to 12\%. A new operation now ships as a kernel module, and the existing verifier analysis stays untouched.